\documentclass[12pt]{article}

% ---------- Packages ----------
\usepackage[margin=1in]{geometry}
\usepackage{graphicx}
\usepackage{fancyhdr}
\usepackage{amsmath}
\usepackage{amssymb}
\usepackage{siunitx}
\usepackage{enumitem}
\usepackage{hyperref}

% ---------- Header/Footer ----------
\pagestyle{fancy}
\fancyhf{}
\fancyhead[L]{Project: 1 Thevenin's Theorem}
\fancyhead[R]{\thepage}
\renewcommand{\headrulewidth}{0pt}

\setlength{\parskip}{0.6em}
\setlength{\parindent}{0pt}

\begin{document}

% ======================================================
% Title Page
% ======================================================
\begin{titlepage}
\thispagestyle{fancy}
\begin{center}
    \vspace*{2cm}
    {\LARGE Linear Circuits 2 \par}
    \vspace{1cm}
    {\Large Thevenin's Theorem Lab Report \par}

    \vspace{6cm}
    {\large Group Members: \par}
    \vspace{1cm}
    {\large Name: Juvon Mondesir \par}
    \vspace{1cm}
    {\large Name: Kenny Gallmon \par}
\end{center}
\end{titlepage}

% ======================================================
% Objectives
% ======================================================
\section*{Objectives:}

The objective of the project was to apply our knowledge of Thevenin and Norton
equivalent circuit that was learnt in Linear Circuits 1, and verify those
results in real world setting, via both simulation of the circuit in LTSpice
as well as physically proving the circuits correct on a breadboard with the
respective components.

The circuit that was to be verified is shown below:

\begin{center}
    \includegraphics[width=0.5\textwidth]{images/circuit_photo.jpg}
\end{center}

% ======================================================
% Equipment
% ======================================================
\section*{Equipment:}

\begin{itemize}[noitemsep]
    \item Digital Multimeter (DMM)
    \item Rohde \& Schwarz Power Supply Series
    \item 1u Capacitor
    \item One 1k Resistor
    \item Two 2k Resistors
    \item One 3k Resistor
\end{itemize}

% ======================================================
% Hand Calculations
% ======================================================
\section*{Hand Calculations:}

To implement this circuit, we first use LTSpice to create it exactly as was
shown via the picture above. Once it was created, it should appear as the
following:

\begin{center}
    \includegraphics[width=0.75\textwidth]{images/fig1_original_circuit.png}
    \par\textit{Figure 1: Original Circuit}
\end{center}

\subsection*{When the circuit is shorted across Node a and b:}

\begin{center}
    \includegraphics[width=0.75\textwidth]{images/fig2_shorted_circuit.png}
    \par\textit{Figure 2: Shorted Circuit}
\end{center}

In the simulation, we found the Thevenin Voltage which would be the
difference in voltage between Node A \& B with the open circuit. As well as
the Norton Current by measuring the current running through the shorted wire
through Nodes A \& B. After this, we then performed a hand analysis to
confirm values, of which the steps are shown below:

\begin{center}
    \includegraphics[width=0.55\textwidth]{images/fig3_vth_sim.png}
    \par\textit{Figure 3: Voltage Thevenin from Simulation}
\end{center}

\begin{enumerate}[label=\arabic*)]
    \item First, remember Ohm's law we know $V = IR$, and $R = V/I$. So, to
    get an idea of each value, we first use the original circuit simulation
    to measure $V_{th}$. Resulting in 2.2V (shown above).
    \item After this, the next value we can obtain is $R_{th}$, which we can
    do multiple ways. The way we chose to do this is by shorting all the
    independent sources and replacing the open circuit between Node A \& B
    with a test voltage of 1V.
    \item Next, using the simulation or KCL, we measure the current going
    across that voltage source which is equivalent to \SI{238.09}{\micro\ampere}.
    \item Once this is done, simply use Ohm's law to get $R_{th}$ by
    calculating $1\text{V} / \SI{238.09}{\micro\ampere}$. The result being
    4200.09 $\Omega$.
    \item This gives us all the components we need to create the Thevenin
    Equivalent Model, but to convert to Norton it's fairly simple.
    \item To get the current needed, divide $V_{th}$ by $R_{th}$ which
    yields the result $5.24 \times 10^{-4}\,\text{A}$ and put a 4200 $\Omega$
    resistor in parallel.
    \item You also get the same result by shorting the wire between A \& B
    and just measuring/solving the current.
\end{enumerate}

\begin{center}
    \includegraphics[width=0.75\textwidth]{images/fig4_hand_calcs.jpg}
    \par\textit{Figure 4: Hand Calculations}
\end{center}

These results are also shown in the Simulation Results section as well.

After this, we then put the circuit onto the breadboard as shown below and
verify it with a Digital Multimeter.

% ======================================================
% Simulation
% ======================================================
\section*{Simulation:}

For the simulation, we decided to build multiple different circuits to
assist us in finding each of our values.

\begin{center}
    \includegraphics[width=0.85\textwidth]{images/fig5_simulated_circuits.png}
    \par\textit{Figure 5: Simulated Circuits}
\end{center}

Continuing to our simulation results, with the following image shown below
for the original circuit:

\begin{center}
    \includegraphics[width=0.7\textwidth]{images/fig6_original_op_point.png}
    \par\textit{Figure 6: Original Circuit Simulation Results}
\end{center}

\textbf{Shown above, are the results for our first circuit. The important
value to gather from this is the voltage V(a), which is 2.2V. (Vth)}

\begin{center}
    \includegraphics[width=0.45\textwidth]{images/fig7_shorted_op_point.png}
    \par\textit{Figure 7: Shorted Circuit Simulation Results}
\end{center}

Here in this image, are the values obtained from the short circuit version.
If you wanted to find $I_{sc}$, it's the equivalent of the current across R3,
approximately $5 \times 10^{-4}$ A. Plugging both the $V_{th}$ and $I_{sc}$
into Ohm's law will give you $R_{th}$.

% ======================================================
% Experiment
% ======================================================
\section*{Experiment:}

For the experimental results, with the breadboard shown above we measure the
voltage between nodes A and B to get the following voltage:

\begin{center}
    \includegraphics[width=0.4\textwidth]{images/fig8_thevenin_voltage.png}
    \par\textit{Figure 8: Thevenin Voltage Measurement}
\end{center}

For the equivalent resistance, we turn off all the independent sources and
use the probes of the Digital Multimeter to get the following equivalent
resistance between nodes A and B:

\begin{center}
    \includegraphics[width=0.4\textwidth]{images/fig9_thevenin_resistance.png}
    \par\textit{Figure 9: Thevenin Resistance Measurement}
\end{center}

For the current, we broke the circuit to measure the current traveling
across the nodes A and B to get the following value:

\begin{center}
    \includegraphics[width=0.4\textwidth]{images/fig10_norton_current.png}
    \par\textit{Figure 10: Norton Equivalent Current (Isc)}
\end{center}

% ======================================================
% Results
% ======================================================
\section*{Results:}

Both the experimental and simulation results are almost identical to each
other, with the only discrepancy being our Thevenin resistance value. After
further inspection our DMM picture has the result of 4.16M $\Omega$ rather
than 4.16k $\Omega$. This could possibly be the result of a broken circuit
accounting for higher resistance than normal given that our calculations and
simulations both prove that 4.16k $\Omega$ is the correct value. Overall,
the results do align with each other, and a simple mistake is most likely
what caused the difference in resistance values given the digits are the
same, just on a larger scale.

% ======================================================
% Conclusion
% ======================================================
\section*{Conclusion:}

This project has been successful in validating the principles of Thevenin
and Norton's theorems via comparing the hands-on calculations, virtual
simulations, and real-world experimentation. The initial hand analysis
determined that the Thevenin voltage to be 2.2V and the Norton current to be
0.5mA, values whose validation was strengthened by the LTSpice simulation
results of 2.2V and 0.52mA, respectively. In the experimental phase, the
physical circuit yielded a measured voltage of 2.195V and a Norton current
of 0.512mA. Overall, the close accuracy across all three validations
confirms the application of circuit theory, making this project a success.

\end{document}